% !PRINTED none

\begin{center}
\gesperrt{Seventh Section.}

Groups of binary linear substitutions.

\rule{4cm}{0.4pt}

\S.\,48.

\gesperrt{Ternary orthogonal substitutions.}
\end{center}

It is now our task to separate out narrower groups from the
totality of the linear substitutions, and finally to arrive at
finite groups. Such narrower groups, which may still be infinite,
are obtained if one makes the requirement that given homogeneous
functions of the variables shall remain invariant. We shall not,
however, pursue the problem further in this generality, but pass
at once to the consideration of the most important special case.
We shall restrict ourselves to ternary substitutions and require
that a quadratic form of non-vanishing determinant shall remain
invariant. Since, as we have seen earlier (Bd. I, \S.\,57), every
such quadratic form can be transformed by linear transformation
into a sum of squares, we do not restrict the problem further if
we take for this quadratic form the sum of the squares. Such
substitutions are called \gesperrt{orthogonal}.

The substitution

\weq{1}{(y_1, y_2, y_3) = A (x_1, x_2, x_3)}

is therefore orthogonal if the substitution coefficients are so determined that the identity holds:

\weq{2}{y_1^2 + y_2^2 + y_3^2 = x_1^2 + x_2^2 + x_3^2.}

If one substitutes the expressions (1) in (2) and sets the
coefficients of corresponding terms equal to one another, then
%CONT>>
